\documentclass[10pt,letterpaper]{article}
\usepackage{cornell-notes}

\title{Clause 31.12.13.40: Weakly Canonical}
\author{}
\date{}
\setDocTitle{Clause 31.12.13.40: Weakly Canonical}
\setDocSubtitle{C++ 2024}
\setDocOwner{C++ 2024 Cornell Notes}
\setDocDate{\today}
\setCornellCollection{C++ 2024 Cornell Notes}
\setCornellUnitType{Clause}
\setCornellUnitNumber{31.12.13.40}
\setCornellUnitTitle{Weakly Canonical}
\hypersetup{pdftitle={Clause 31.12.13.40: Weakly Canonical},pdfsubject={Cornell notes},pdfauthor={}}

\begin{document}
\maketitle
\begin{CornellOverviewBox}{Overview}
\textbf{Classification:} Standard library - Normative\\[2pt]
\textbf{Learning objective:} Understand the rules, terminology, and semantic consequences associated with Weakly canonical.
\end{CornellOverviewBox}\begin{CornellSummaryBox}{Summary}
{\sffamily\bfseries\color{Accent} Summary}\par\vspace{3pt}Section 31.12.13.40 focuses on Weakly canonical. Apply the specified interface together with its mandates, effects, result conditions, exception behavior, and complexity constraints. When applying it, preserve the distinction between mandatory requirements, recommendations, permissions, and behavior deliberately left undefined, unspecified, or implementation-defined.
\end{CornellSummaryBox}

\end{document}
